\chapter{Princípios e Fundamentos}

O curso de \ppccurso da Universidade Federal de Uberlândia é orientado por princípios que articulam a excelência acadêmica com a responsabilidade social, a inovação com a ética, a técnica com a criticidade. Tais princípios refletem os compromissos institucionais da UFU, as Diretrizes Curriculares Nacionais da Computação e os desafios da formação de engenheiros no século XXI.

Entre os fundamentos que estruturam o presente Projeto Pedagógico do Curso, destacam-se:

\begin{itemize}
    \item Formação científica, técnica e humanística sólida, com domínio teórico e prático das bases da computação e da engenharia, articuladas à capacidade crítica de compreender e intervir no mundo de forma ética e responsável;
    \item Integração entre hardware e software, com abordagem sistêmica e interdisciplinar, valorizando múltiplos paradigmas computacionais e suas aplicações em contextos diversos;
    \item Articulação entre ensino, pesquisa e extensão, como indissociabilidade fundadora da universidade pública, promovendo a formação integral do(a) estudante em contato com problemas reais e demandas da sociedade;
    \item Compromisso com a inovação, com incentivo à criatividade, à autonomia intelectual, à cultura empreendedora e ao uso de tecnologias emergentes como instrumentos de transformação social;
    \item Inclusão e diversidade, compreendidas como dimensões fundamentais da formação acadêmica, que exigem o acolhimento de múltiplas trajetórias, saberes e identidades, bem como a construção de ambientes educativos plurais, acessíveis e antidiscriminatórios;
    \item Adoção de metodologias ativas e educação híbrida, combinando estratégias presenciais e digitais de ensino-aprendizagem, promovendo a flexibilidade curricular, a personalização da formação e a ampliação do repertório pedagógico;
    \item Formação orientada por competências -- \textit{Competency-Based Learning} (CBL) --, com foco em resultados de aprendizagem significativos, articulação entre saberes, habilidades e atitudes, e ênfase na capacidade de mobilizar conhecimentos para resolver problemas complexos em contextos reais;
    \item Formação orientada pela sustentabilidade, ética e responsabilidade social, com foco na resolução de problemas complexos e multifacetados, em alinhamento com os princípios do \textit{Engineering Education for Sustainable Development} (EESD) e dos Objetivos de Desenvolvimento Sustentável (ODS)\footnote{Disponível em: \url{https://brasil.un.org/pt-br/sdgs}}, capacitando o(a) estudante a integrar dimensões técnicas, sociais, ambientais e econômicas em sua atuação profissional e cidadã.
\end{itemize}

Esses princípios orientam as decisões curriculares, metodológicas e avaliativas do curso, guiando a formação de engenheiros e engenheiras de computação capazes de atuar de forma crítica, criativa e transformadora nos diversos contextos em que a tecnologia se insere.

\section{A UFU e o Curso}

A Universidade Federal de Uberlândia (UFU) estabelece, como fundamentos institucionais, a indissociabilidade entre ensino, pesquisa e extensão, o compromisso com a formação cidadã e crítica, e a valorização do conhecimento como instrumento de transformação social. Sua missão institucional é definida nos seguintes termos:

\begin{displayquote}
\textit{Desenvolver o ensino, a pesquisa e a extensão de forma integrada, realizando a função de produzir e disseminar as ciências, as tecnologias, as inovações, as culturas e as artes, e de formar cidadãos críticos e comprometidos com a ética, a democracia e a transformação social.}
\end{displayquote}

A visão que orienta sua atuação projeta a instituição como:

\begin{displayquote}
\textit{Ser referência regional, nacional e internacional de universidade pública na promoção do ensino, da pesquisa e da extensão em todos os campi, comprometida com a garantia dos direitos fundamentais e com o desenvolvimento regional.}
\end{displayquote}

O curso de \ppccurso compartilha dessa missão e contribui ativamente para sua realização. Sua proposta pedagógica reflete os valores centrais da UFU, como o compromisso com a excelência acadêmica, a equidade, a responsabilidade social, a inovação e a formação de profissionais comprometidos com o desenvolvimento sustentável e com a transformação da realidade em que estão inseridos.

A construção e reformulação do curso têm sido conduzidas de forma democrática e participativa, envolvendo docentes, técnicos-administrativos, discentes e representantes da comunidade acadêmica e regional. As decisões curriculares e metodológicas refletem essa escuta ampliada e estão em consonância com os planos institucionais da universidade, como o Plano Institucional de Desenvolvimento e Expansão (PIDE) e as diretrizes do CONSUN e do CONGRAD.

Há, portanto, plena coerência entre os princípios do curso e os ideais da UFU. A formação ofertada pela \ppccurso articula conhecimentos técnico-científicos com o compromisso ético e social, contribuindo para a consolidação da universidade como espaço de produção de conhecimento crítico, emancipador e transformador.